This work investigates how explainable artificial intelligence methods
can be applied for medical image classification rather than used only as
visual illustrations. The study focuses primarily on chest X-ray
pneumothorax detection using the SIIM-ACR pneumothorax dataset and
off-the-shelf TorchXRayVision classifiers, with a completed head CT
intracranial hemorrhage pilot used as a cross-modality methodological
extension. The aim is to compare several post-hoc explanation methods
against out-of-the-box medical classification models and assess whether
their highlighted evidence is localized, faithful to the model, and
clinically useful.

The methodology separates classifier behavior from explanation quality.
Pneumothorax classifiers are evaluated at frozen thresholds, while
explanation maps are generated with CAM-family, SHAP and Occlusion
methods, and a frozen four-method consensus. Each method is interpreted
through positive, negative, magnitude, and signed views. Explanations
are validated with lesion-mask localization metrics, negative-evidence
diagnostics, deletion/insertion faithfulness curves, method-agreement
analysis, and a structured radiologist-style review of selected cases.

Current CXR results show that DenseNet-chex supersedes the original
DenseNet-all checkpoint as the selected DenseNet baseline, while
ResNet-50 remains the strongest tested TorchXRayVision follow-up in
relative localization terms; nevertheless, absolute pneumothorax-mask
overlap remains weak. The held-out Phase 5.2 improvement experiment
shows model-dependent consensus behavior: DenseNet-chex consensus mainly
improves over Score-CAM in paired overlap tests, whereas ResNet-50
consensus significantly improves Dice/IoU over several weaker individual
methods but not over the strongest CAM comparators. The balanced 40-case
review shows mixed usefulness: many maps contain interpretable evidence,
yet clinically misleading, indirect, device-related, or high-contrast
non-lesion evidence remains common.

The practical significance of the work is a reproducible validation
workflow for auditing medical image explanations. The thesis argues that
heatmaps should be treated as model-behavior diagnostics, not anatomical
segmentations or automatic generators of clinical trust.

Keywords: explainable artificial intelligence, medical image
classification, radiology, Grad-CAM, SHAP, pneumothorax, intracranial
hemorrhage
